\subsection{The full relative pro-\texorpdfstring{\(p\)}{p} cover over an odd finite base}

\begin{statement}{D.41}{Proved}{The exact pairing-compatible abelian image and all relative pro-$p$ peripheral lifts}
Let \(q:\widehat F\twoheadrightarrow D\) have finite odd-order
image, let \(K=\ker q\), and let \(p\nmid |D|\).
Write \(R=K^{(p)}\) for the maximal pro-\(p\) quotient of the
actual open subgroup \(K\), and put
\[
 E_p=\widehat F/\ker(K\to R),\qquad
 E_{p,\mathrm{ab}}=E_p/\overline{[R,R]}.
 \tag{D.41.1}
\]
Then \(R\) is free pro-\(p\) of rank \(|D|+1\), and
\(1\to R\to E_p\to D\to1\) is the actual source extension.

Fix \(\alpha\in\operatorname{Aut}(D)\) and specified
\(h_i\in D\) with
\(\alpha(d_i)=h_i d_i h_i^{-1}\), and use multiplier one.
Every pairing-compatible automorphism of
\(E_{p,\mathrm{ab}}\) above \(\alpha\), with these specified
peripheral conjugators, lifts to a peripheral automorphism
of \(E_p\) above the same \(\alpha\) and with the same
conjugator projections.
Here pairing-compatible means that its induced map on
the compact homology \(C_p\) of D.40 preserves the
actual integral intersection form.
Conversely every such automorphism of \(E_p\) has this
property.  This gives the exact image of the full relative
pro-\(p\) peripheral fibre in its abelian fibre.

The fibre is nonempty for every specified
\((\alpha,h_1,h_2,h_3)\).
Every generating source map
\(\widehat F\to E_p\) above \(\alpha q\) factors through
an automorphism of this same \(E_p\).
Its full literal peripheral fibre has the exact
normalizer parameterization in
\textup{(D.41.22)--(D.41.24)} below.
In particular all D.37 weight-five shadows admit
compatible generating lifts to the full relative
pro-\(p\) cover for every \(p\ne5\), with their
specified characteristic five-group quotient retained.
\end{statement}

\begin{proof}
Put \(n=|D|\), and let \(b=\sum_i n/e_i\) be the number
of boundary components in the regular pants cover.
There are at least three such components.
Finite-coset Schreier rewriting identifies \(K\)
with the free profinite group of rank \(n+1\)
and identifies its maximal pro-\(p\) quotient
with the free pro-\(p\) group \(R\) on the same
Schreier generators.
The kernel of \(K\to R\) is characteristic in \(K\),
hence normal in \(\widehat F\).
This proves the exact extension in
\textup{(D.41.1)}, with the indicated source map.
Its boundary generators are the literal based words
\[
 c_\gamma=t_d x_i^{e_i}t_d^{-1}\in R,
       \qquad
 \gamma=d\langle d_i\rangle.
 \tag{D.41.2}
\]
Their conjugacy classes are independent of the
chosen coset paths, and their homology classes
are the individually labelled boundary vectors
of D.39.14.

\paragraph{The actual boundary automorphism group and its integral image.}
Let \(\mathcal H\) be the subgroup of
\(\operatorname{Aut}(R)\) whose elements permute
the conjugacy classes of the \(c_\gamma\), retaining
their positive generators, rather than replacing
them by unspecified unit powers.
The group \(\operatorname{Aut}(R)\) is compact:
its image in \(R^{n+1}\) consists exactly of the
tuples whose reductions form a basis of
\(R/\Phi_p(R)\), a clopen condition.
Such a tuple generates \(R\); a proper closed
subgroup would be contained in an open maximal
subgroup, contradicting that basis condition.
The resulting surjective endomorphism is an
automorphism, as follows from the characteristic
finite \(p\)-central quotients below.
Each conjugacy class of \(c_\gamma\) is a compact
image of \(R\).  Thus each prescribed permutation
condition is closed, and their finite union
\(\mathcal H\) is a closed profinite subgroup.

The integral homology is \(M_p\), with boundary
radical \(B_p\) and compact quotient \(C_p\).
Let \(\mathcal G(\mathbf Z_p)\) be the group of
linear automorphisms of \(M_p\) that permute the
individual boundary vectors and induce an
isometry of \(C_p\).
We claim the exact equality
\[
 \operatorname{im}\bigl(\mathcal H\to
               \operatorname{Aut}_{\mathbf Z_p}(M_p)\bigr)
                       =\mathcal G(\mathbf Z_p).
 \tag{D.41.3}
\]
We prove both inclusions with actual maps.

For necessity, take two homology characters
vanishing on the boundary vectors, with values
in \(\mathbf Z/p^a\).
Send each free generator of \(R\) to
\((\chi_j,\psi_j,0)\) in the finite Heisenberg
\(p\)-group.  Every boundary image is central.
In a surface presentation the sum of their
central values is, up to the one common
orientation sign, the handle cup pairing.
Changing the chosen Heisenberg lift of
\((\chi,\psi)\) adds a linear character; its
sum on the boundaries is zero by their
single homology relation.
An element of \(\mathcal H\) permutes those
central boundary values and conjugates them,
which does not change their sum.
It therefore preserves the cup form modulo
every \(p^a\), and hence preserves the
integral intersection form on \(C_p\).
This proves the inclusion in
\textup{(D.41.3)}.

For the reverse inclusion use the actual
oriented surface presentation
\[
 \prod_{j=1}^g[a_j,b_j]\,c_1\cdots c_b=1.
 \tag{D.41.4}
\]
Omitting \(c_b\) gives a free basis of \(R\).
Adjacent boundary permutations are the literal
Hurwitz automorphisms
\[
 (c_j,c_{j+1})\longmapsto
    (c_{j+1},c_{j+1}^{-1}c_jc_{j+1}),
 \tag{D.41.5}
\]
with inverse
\((c_j,c_{j+1})\mapsto
(c_jc_{j+1}c_j^{-1},c_j)\).
They preserve the product in
\textup{(D.41.4)} exactly, including when one
boundary is the omitted generator.
They realize every permutation of the
oriented boundary classes.

On the handle part, a Dehn twist is the
explicit homeomorphism obtained by cutting
an annular neighbourhood of an oriented
simple curve, rotating one boundary once,
and regluing, with the surface base point
outside that annulus.
Its induced free-group automorphism extends
to \(R\), and its homology map is the
transvection
\(v\mapsto v+\langle v,a\rangle a\), with
the sign chosen by the direction of the twist.
The curves for the standard \(a_j,b_j\),
and band sums in two handles, realize the
primitive vectors \(a_i,b_i,a_i+a_j,a_i+b_j\)
and \(b_i+b_j\), with \(i\ne j\) in every
cross-handle sum.
The two twists within a handle can also be
written literally as
\[
 (a,b)\mapsto(a,ba),\qquad
 (a,b)\mapsto(ab^{-1},b);
 \tag{D.41.6}
\]
both preserve \([a,b]\).
These elementary symplectic transvections
generate \(\operatorname{Sp}_{2g}(\mathbf Z/p^a)\)
at every \(a\).
Indeed a unimodular first column over the
local ring has a unit entry; symplectic swaps,
the two elementary operations on a handle,
and the indicated cross-handle transvections
move that entry to the first handle, scale it
to one, and clear the other entries.
The paired column is then cleared using the
symplectic pairing, and the same operations
apply to the remaining \(g-1\) handles.
For complete explicitness, write
\(E_{12}(t)=\left(\begin{smallmatrix}1&t\\0&1\end{smallmatrix}\right)\)
and
\(E_{21}(t)=\left(\begin{smallmatrix}1&0\\t&1\end{smallmatrix}\right)\).
For every unit \(u\) in the local ring,
\[
 \begin{pmatrix}u&0\\0&u^{-1}\end{pmatrix}
 =
 E_{12}(u-1)E_{21}(1)
 E_{12}(u^{-1}-1)E_{21}(-u).
\]
Direct multiplication proves this identity,
so the unit scaling also lifts by integer
twist powers.  This is an induction on \(g\).
Consequently these actual integral maps have
dense image in \(\operatorname{Sp}(C_p)\).

We retain explicit maps for the remaining
\(\operatorname{Hom}(C_p,B_p)\) part.
In a three-generator subsurface word
\([a,b]c\), the automorphism
\[
 \begin{split}
 a&\longmapsto c^{-1}ac,\\
 b&\longmapsto c^{-1}a^{-1}cab c,\\
 c&\longmapsto c^{-1}a^{-1}ca c
 \end{split}
 \tag{D.41.7}
\]
preserves \([a,b]c\) exactly and sends
\(c\) to a positive conjugate of itself.
Here is an actual inverse construction.
Use the free basis
\[
 u=ac,\qquad v=c^{-1}bc,\qquad z=vcv^{-1},
 \tag{D.41.8}
\]
whose inverse is
\[
 c=v^{-1}zv,\qquad
 a=uv^{-1}z^{-1}v,\qquad
 b=v^{-1}zvz^{-1}v.
 \tag{D.41.9}
\]
In that basis \([a,b]c=[u,v]z\).
Conjugate the twist \(v\mapsto vu\) through
\textup{(D.41.8)--(D.41.9)}, and compose it
with the inverse of \(b\mapsto ba\).
This gives \textup{(D.41.7)}.
Replacing the two twists by their inverse
maps in reverse order gives its inverse.
Its homology action is precisely
\[
                   b\longmapsto b+c,\qquad
                   a\longmapsto a,\qquad
                   c\longmapsto c.
 \tag{D.41.10}
\]
For a boundary \(c_j\) separated from the
first handle in \textup{(D.41.4)} by a word
\(W\) in the other generators, replace that
boundary temporarily by \(Wc_jW^{-1}\).
This is an invertible basis replacement
and puts the local product in the form
\([a,b](Wc_jW^{-1})W\).
Apply \textup{(D.41.7)} and then undo the
basis replacement.  It preserves the full
surface relation and all positive boundary
classes, with homology addition \(b\mapsto b+c_j\).
Conjugating by the handle maps and taking
integer powers gives every elementary
map from the compact handle lattice to
the independent \(b-1\) boundary vectors.
These maps add because each vanishes on
the boundary lattice.  Their closure is
the full group \(\operatorname{Hom}_{\mathbf Z_p}(C_p,B_p)\).

In the surface homology basis, an element
of \(\mathcal G(\mathbf Z_p)\) is determined
by its boundary permutation, its symplectic
compact block, and this arbitrary off-diagonal
homomorphism.
The preceding maps therefore have dense
image in \(\mathcal G(\mathbf Z_p)\).
The image of compact \(\mathcal H\) is
closed, proving equality
\textup{(D.41.3)}.

\paragraph{The exact pro-\(p\) kernel and coprime conjugacy.}
The kernel of
\(\operatorname{Aut}(R)\to
\operatorname{GL}(R/\Phi_p(R))\)
is pro-\(p\).
Here is a finite argument that also proves
the needed Hopfian assertion.
Use the characteristic lower \(p\)-central
series \(R_1=R\),
\(R_{j+1}=\overline{R_j^p[R_j,R]}\).
Its finite quotients are cofinal among
finite \(p\)-quotients: the analogous series
terminates in every finite \(p\)-group.
For a finite \(p\)-group \(U\) on \(r\)
generators, the automorphisms inducing the
identity on \(U/\Phi_p(U)\) act freely on
the ordered generating tuples lifting a
fixed Frattini basis.
There are \(|\Phi_p(U)|^r\) such tuples,
so the order of this automorphism kernel
divides a power of \(p\).
Applying this to \(R/R_j\) proves the
pro-\(p\) assertion.
A surjective endomorphism of \(R\) induces
automorphisms of all those characteristic
finite quotients; their cofinality makes
its kernel trivial.
In particular
\[
 \mathcal P=
 \ker\bigl(\mathcal H\to
              \mathcal G(\mathbf Z_p)\bigr)
                 \quad\text{is pro-}p.
 \tag{D.41.11}
\]

We use the following exact coprime fact.
In an extension with normal pro-\(p\)
kernel and finite quotient of order prime
to \(p\), complements exist and any two
are conjugate by the kernel.
For a finite \(p\)-kernel, prove it
inductively on its order using the nonzero
characteristic elementary abelian group
\(\Omega_1 Z(P)\).
After the induction step in the quotient,
the remaining existence and conjugacy
equations are respectively \(H^2(D,A)=0\)
and \(H^1(D,A)=0\) for that elementary
abelian module \(A\).
They have explicit averaging solutions:
for a normalized two-cocycle \(f\), put
\(b(d)=|D|^{-1}\sum_e f(d,e)\);
then \(f(d,e)=b(d)+d b(e)-b(de)\).
For a crossed homomorphism \(z\), put
\(v=|D|^{-1}\sum_e z(e)\);
then \(z(d)=v-dv\).
These formulas split the extension and
conjugate its two complements.
For a pro-\(p\) kernel, apply the finite
argument in all finite quotients and use
compactness: the closed sets of compatible
section lifts, or of conjugating elements
in the kernel, have the finite intersection
property.  Their intersections give the
required actual complement and conjugator.

Choose such a complement in \(E_p\).
Write its section as \(s:D\to E_p\).
It identifies the actual extension with
\[
                    E_p=R\rtimes_\rho D.
 \tag{D.41.12}
\]
The action \(\rho(D)\) lies in
\(\mathcal H\), since conjugation by
source lifts permutes the boundary
classes of the regular cover.
It is faithful already on homology
modulo \(p\), by D.39.18, which contains
the regular \(D\)-module.

\paragraph{Making a prescribed homology lift exactly equivariant.}
Take any prescribed pairing-compatible
automorphism of \(E_{p,\mathrm{ab}}\)
in the statement.
Normalize its action on the finite
complement by an inner element of \(M_p\):
the averaging formula for \(H^1(D,M_p)\)
used in D.39.23 gives this explicitly.
Its homology restriction is then an
\(\alpha\)-semilinear map
\[
 T\in\mathcal G(\mathbf Z_p),\qquad
 T\rho_M(d)T^{-1}=\rho_M(\alpha(d)),
 \tag{D.41.13}
\]
with the exact boundary permutation
\[
 d\langle d_i\rangle\longmapsto
                \alpha(d)h_i\langle d_i\rangle.
 \tag{D.41.14}
\]
By \textup{(D.41.3)} choose an actual
\(\sigma\in\mathcal H\) inducing \(T\).
In the preimage of the finite subgroup
\(\rho_M(D)\) under
\(\mathcal H\to\mathcal G(\mathbf Z_p)\),
both \(\rho(D)\) and
\(\sigma\rho(D)\sigma^{-1}\)
are complements to the pro-\(p\)
kernel \(\mathcal P\).
The coprime conjugacy just proved
gives \(u\in\mathcal P\) such that
\(\eta=u\sigma\) satisfies
\[
                  \eta\rho(d)\eta^{-1}
                         =\rho(\alpha(d))
                           \quad(d\in D).
 \tag{D.41.15}
\]
Equality with these labels follows
from \textup{(D.41.13)} and faithfulness
of the finite homology image.
The correction \(u\) acts trivially on
all integral homology, so \(\eta\)
still induces the prescribed \(T\).
It also retains every boundary label:
distinct boundary generators have
distinct homology vectors, by the
single sum relation and \(b\geq3\).
Consequently
\[
                   (r,d)\longmapsto
                        (\eta(r),\alpha(d))
 \tag{D.41.16}
\]
is an actual automorphism of \(E_p\).

\paragraph{Recovering the peripheral generators from their powers.}
This step retains the exact conjugators.
Let \(b_i=X_i^{e_i}\in R\) and take
the specified lift \(H_i\in E_p\) of
\(h_i\).
By the positive boundary-class property
of \(\eta\) and \textup{(D.41.14)},
there is \(a_i\in R\) with
\[
       \eta(b_i)=a_iH_i b_iH_i^{-1}a_i^{-1}.
 \tag{D.41.17}
\]
Set \(Y_i=\widetilde\alpha(X_i)\) and
\(U_i=a_iH_i X_iH_i^{-1}a_i^{-1}\).
They have the same image in \(D\) and
the same \(e_i\)th power
\[
                        Y_i^{e_i}=U_i^{e_i}=v_i.
 \tag{D.41.18}
\]

Every boundary generator is primitive in
\(R\).  This is immediate for the
\(b-1\) independent boundaries in
\textup{(D.41.4)}; the last is primitive
as well, because another boundary
occurs exactly once in its surface word.
Replacing that other boundary by the
last one is a free basis replacement
with inverse obtained from
\textup{(D.41.4)}.
For a primitive boundary \(c\),
\[
                            C_R(c)=\overline{\langle c\rangle}.
 \tag{D.41.19}
\]
To verify the exact centralizer, take
\(c\) as the first free pro-\(p\)
generator and send the generators
to \(1+X_j\) in
\(\mathbf F_p\langle\!\langle X_1,\ldots,X_{n+1}\rangle\!\rangle^\times\).
This Magnus map is injective.
Indeed any finite \(p\)-group quotient
can be evaluated by substituting
\(X_j\mapsto g_j-1\) in its group algebra;
the augmentation ideal of a finite
\(p\)-group algebra is nilpotent.
The latter follows by induction using
a central element of order \(p\),
whose difference from one has \(p\)th
power zero.
Thus equality in the Magnus algebra
implies equality in every finite
\(p\)-quotient, and hence in \(R\).
Comparing coefficients in
\(X_1A=AX_1\) shows that the centralizer
of \(X_1\) in the power-series algebra
is exactly \(\mathbf F_p[[X_1]]\):
the coefficient of a word containing
another letter is shifted to one
having no initial \(X_1\), where the
commutation identity forces zero.
The retraction \(R\to\langle c\rangle\)
sets the other \(X_j\) to zero.
For a commuting group element it
leaves the Magnus series unchanged,
and injectivity proves
\textup{(D.41.19)}.

Now \(Y_iU_i^{-1}\) lies in \(R\)
and centralizes \(v_i\).
The element \(v_i\) is a conjugate
of a primitive boundary generator,
so this difference belongs to
the procyclic group \(\langle v_i\rangle\).
Both \(Y_i\) and \(U_i\) fix \(v_i\)
and hence commute with that entire
procyclic group.
Taking \(e_i\)th powers in
\textup{(D.41.18)} therefore gives
\((Y_iU_i^{-1})^{e_i}=1\).
The group \(\langle v_i\rangle\simeq
\mathbf Z_p\) is torsion-free.
It follows that
\[
                 \widetilde\alpha(X_i)
                    =a_iH_iX_iH_i^{-1}a_i^{-1}.
 \tag{D.41.20}
\]
Thus the actual peripheral roots,
and their specified base conjugators,
are retained, rather than only their
boundary powers.

Finally lift the abelian normalizing
element used before \textup{(D.41.13)}
from \(M_p\) to \(R\), and conjugate
\textup{(D.41.16)} by it.
This gives a lift of the entire
prescribed abelian automorphism,
including its section-conjugator
coordinate.
Necessity of the intersection condition
follows from \textup{(D.41.3)}.
D.40 supplies a pairing-compatible
abelian lift for every specified
\((\alpha,h_i)\), proving nonemptiness.

\paragraph{All source maps and the exact full fibre.}
Any source homomorphism
\(f:\widehat F\to E_p\) above
\(\alpha q\) maps \(K\) into \(R\).
By the universal property of
\(K^{(p)}\), it kills \(\ker(K\to R)\)
and factors through the exact quotient
\textup{(D.41.1)}.
If \(f\) is onto, its restriction
on the kernel \(R\) is onto: a
preimage of an element in \(R\)
must have trivial image in \(D\),
since the base map is \(\alpha\).
The Hopfian argument above makes
this restriction an automorphism,
and then the whole induced map on
\(E_p\) is an automorphism.
The identical argument on the
abelian kernel \(M_p\) applies to
source epimorphisms to
\(E_{p,\mathrm{ab}}\).

Here is the literal full-fibre
parameterization, including its
remaining conjugator quotient.
Let \(\mathcal H_0\) be the subgroup
of \(\mathcal H\) fixing each positive
boundary conjugacy class, and let
\(\theta_h\) be the permutation
\textup{(D.41.14)}.
Define the closed set
\[
 \mathcal Y_{\alpha,h}=
 \{\eta\in\mathcal H:
   \eta\rho(d)\eta^{-1}=\rho(\alpha(d))
       \text{ for all }d,\quad
   \eta\text{ induces }\theta_h\}.
 \tag{D.41.21}
\]
It is nonempty by the construction.
For any one of its elements \(\eta_0\),
\[
 \mathcal Y_{\alpha,h}
      =C_{\mathcal H_0}(\rho(D))\,\eta_0.
 \tag{D.41.22}
\]
This equality follows by multiplying
the two exact intertwining identities;
the boundary permutations cancel.
Every automorphism in the full
fixed-conjugator fibre is
\[
      \operatorname{Inn}(r)\circ
          \bigl((s,d)\mapsto(\eta(s),\alpha(d))\bigr),
         \qquad r\in R,\quad
                \eta\in\mathcal Y_{\alpha,h}.
 \tag{D.41.23}
\]
Indeed its image of the chosen
finite complement is another
complement, so coprime conjugacy
normalizes it to \(\rho(D)\).
The ambiguity is exactly
\[
 (r,\eta)\sim
       (ra,\operatorname{Inn}(a^{-1})\eta),
       \qquad
       a\in R^D=C_R(s(D)).
 \tag{D.41.24}
\]
These pairs give the same automorphism,
and equality on the finite complement
and on \(R\) proves the converse.
This is a quotient of compact spaces
by the displayed actual group action.

Allowing ordinary base conjugators
\(h_ic_i\), \(c_i\in C_D(d_i)\),
gives the disjoint union over the
distinct boundary permutations;
their number is the product in
D.39.6, by the same individually
distinguishable boundary vectors.
This describes the entire source
fibre at the relative pro-\(p\)
stage, not a selected independently
admitted component.
When \(D\) is a five-group and
\(p\ne5\), the projection \(E_p\to D\)
is its maximal pro-five quotient:
the normal pro-\(p\) group \(R\)
has trivial image in any finite
five-group.  This proves the exact
fixed-target application to D.37.
\end{proof}

The finite program
\texttt{a\_relative\_prop.py}
checks \textup{(D.41.5)--(D.41.10)}
as literal free-word maps, with
their inverses, including a swap
involving the omitted boundary.
It generates all 720 matrices
of \(\operatorname{Sp}_4(\mathbf F_2)\)
and all 24 matrices of
\(\operatorname{Sp}_2(\mathbf F_3)\)
from exactly the stated elementary
transvections: eight listed vectors in
genus two over \(\mathbf F_2\) and the
two handle vectors in genus one over
\(\mathbf F_3\).  The complete vector
and matrix lists are retained in its JSON
certificate.
It also checks all 432
automorphisms of the actual
Heisenberg group over \(\mathbf F_3\),
whose identity-on-Frattini kernel
has order nine, and exhibits the
four kernel conjugators between
the four complements in
\(\mathbf F_2^2\rtimes C_3\).
The infinite pro-\(p\) assertions
use the proved compactness and
characteristic-quotient arguments,
rather than extrapolation from
these finite checks.

D.41 controls one entire relative
pro-\(p\) kernel over the specified
odd-order base.  A finite quotient
whose further kernel mixes this
prime with additional chief factors
is not supplied by the maps
\textup{(D.41.1)}.
The remaining D.1--D.2 comparison
must retain those further finite
quotients and the complete fibres
over the pairing-compatible
relative lifts just constructed.
